% !TeX root = ../../Book2.tex
\chapter{分离变量与 Sturm–Liouville 理论}
\label{b2:ch:10}
\BookChapterNumber{b2:ch:10}
% 本章摘要（不计入编号）
\begin{chapterabstract}
分离变量把演化方程中的空间形状与时间变化分开，但一般初值能否如此展开，需要独立证明。本章先计算有限区间上三类边界条件的模态，再在连续正系数及分离端点条件下构造 Sturm--Liouville 算子的移位 Green 核。紧自伴谱定理给出离散谱和完备基，能量空间的系数刻画随后把形式级数提升为波动与热传导问题的范数极限解。推导中将分别说明边界符号、零模态和负特征值的影响。
\end{chapterabstract}

% 本章目标（不计入编号）
\begin{learninggoals}
\begin{itemize}
\item\label{b2:item:10-goal-modes}计算 Dirichlet、Neumann、Robin 及混合边界的空间模态，并判断零频率是否存在。
\item\label{b2:item:10-goal-domain}区分形式对称性与无界算子的自伴性，写清加权空间、定义域和分离条件。
\item\label{b2:item:10-goal-green}从两个常微分方程基本解构造移位 Green 核，解释紧性和单射性如何给出完备性。
\item\label{b2:item:10-goal-energy}用系数加权平方和分别刻画算子定义域与能量空间。
\item\label{b2:item:10-goal-limit}通过有限解的稳定估计构造广义解，说明范数极限中保存了哪些初边值信息。
\end{itemize}
\end{learninggoals}

本章使用\chapref{b2:ch:02}的正交展开与\secref{b2:sec:0104}中初边值问题的定义。紧算子部分以前卷的紧自伴谱定理为前置，具体为第一卷定理 8.5.5；一维 Sobolev 代表取自第一卷命题 21.4.1。未读过第一卷选读谱理论的读者，应先补读该定理。阅读本章不需要预先掌握一般无界谱定理或后编的变分存在性理论。

\ProseParagraphBreak

\secref{b2:sec:1001}从常系数模型讲起，\secref{b2:sec:1002}构造移位逆并确定算子的自伴定义域，\secref{b2:sec:1003}证明特征函数的完备性，\secref{b2:sec:1004}研究相应的演化级数及其极限。John K. Hunter 与 Bruno Nachtergaele 的《Applied Analysis》\cite[Section 10.5, pp. 268--272]{Hunter2001Applied}讨论 Green 逆与特征展开，先取 \(p\in C^1\)、权重为一的 Dirichlet 问题。本章将通量作为变量，明确处理连续 \(p\)、正权重和实 Robin 系数。分离变量法在各类边界模型中的应用，也可在 Walter A. Strauss 的《Partial Differential Equations: An Introduction》\cite{Strauss2008Partial}中阅读。

% !TeX root = ../../Book2.tex
\section{分离变量法}
\label{b2:sec:1001}

\subsection{有限区间}
\label{b2:subsec:100101}

同一根细杆上的振动与导热具有不同的时间规律，却受同一种空间边界约束控制。设杆占据区间 \(I=(0,\ell)\)，\(\ell>0\)。常系数模型分别是 \(u_{tt}=c^2u_{xx}\) 与 \(u_t=\kappa u_{xx}\)，其中 \(c,\kappa>0\)。若先寻找非零乘积 \(u(t,x)=T(t)X(x)\)，在两因子不为零的小矩形上相除，便得到一个只依赖 \(t\)、又只依赖 \(x\) 的量。它必须是常数；把空间方程写成
\[
 -X''=\lambda X,
 \quad T''+c^2\lambda T=0
 \quad\text{或}\quad T'+\kappa\lambda T=0.
\]
这称为\term[fenli-bianliang-fa]{分离变量法}[separation of variables]。实际计算不必在零点反复相除：一旦得到这两个常微分方程，将乘积直接代回原方程，就能在整个区域验证它。

\ProseParagraphBreak

分离变量首先产生某些特殊解。它既没有断言每个解都是一个乘积，也没有保证任意初值都能展开为这些空间因子。前一个问题由线性叠加处理，后一个问题需要本章的完备性定理。有限和一定满足同一齐次线性方程；无穷和是否仍满足方程，则还取决于收敛发生在哪个空间。

\ProseParagraphBreak

齐次边界条件应先施加于 \(X\)。如果端点温度给定为非零常数，一般不能让每个分离项各自满足它们，而应先减去一个满足边界值的函数。相应的非齐次问题将在\chapref{b2:ch:12}中处理。下面先把三种齐次边界条件对应的频率区分清楚。

\subsection{Dirichlet 边界条件}
\label{b2:subsec:100102}

两端固定或温度保持为零时，空间条件为 \(X(0)=X(\ell)=0\)。对特征方程乘 \(\overline X\) 并积分，得到
\[
 \lambda\int_0^\ell |X|^2\,\mathrm{d}x
 =\int_0^\ell |X'|^2\,\mathrm{d}x.
\]
非零解因此只能对应实数 \(\lambda>0\)：若 \(\lambda=0\)，导数为零，端点条件再迫使常数为零。写 \(\lambda=k^2\)，通解是 \(A\cos(kx)+B\sin(kx)\)。左端条件给出 \(A=0\)，右端条件给出 \(k\ell=n\pi\)，\(n\geq1\)。归一化后为
\begin{equation}\label{b2:eq:10-dirichlet-modes}
 e_n(x)=\sqrt{\dfrac2\ell}\sin\dfrac{n\pi x}{\ell},
 \quad\lambda_n=\left(\dfrac{n\pi}{\ell}\right)^2,
 \quad n\geq1.
\end{equation}
直接积分验证 \(\langle e_n,e_m\rangle=\delta_{nm}\)。这里的正交性还只是两两关系，完备性将在\secref{b2:sec:1003}中作为一般定理的特例重新取得。

\ProseParagraphBreak

一个模态在振动时按 \(\cos(c\sqrt{\lambda_n}t)\) 或 \(\sin(c\sqrt{\lambda_n}t)\) 变化，在导热时按 \(e^{-\kappa\lambda_nt}\) 衰减。因此同一空间高频在前一个方程中振动更快，在后一个方程中衰减更快。区间放大 \(r\) 倍，特征值缩小为原来的 \(r^{-2}\)；振动时间尺度随长度变化，扩散时间尺度则随长度平方变化。

\subsection{Neumann 边界条件}
\label{b2:subsec:100103}

绝热端点对应 \(X'(0)=X'(\ell)=0\)。分部积分仍给出非负特征值，但这次常数不再被排除。零特征值有单位特征函数 \(e_0=\ell^{-1/2}\)，而正特征值及单位特征函数为
\begin{equation}\label{b2:eq:10-neumann-modes}
 \lambda_n=\left(\dfrac{n\pi}{\ell}\right)^2,
 \quad e_n(x)=\sqrt{\dfrac2\ell}\cos\dfrac{n\pi x}{\ell},
 \quad n\geq1.
\end{equation}
求解时必须单独保留 \(n=0\)，不能把含 \(\lambda_n^{-1}\) 或 \(\lambda_n^{-1/2}\) 的正频公式直接用于它。热方程的零模态保持不变；波方程的零模态为 \(a_0+b_0t\)，代表整体平移而没有恢复力。

\ProseParagraphBreak

左右端也可以采用不同类型。例如 \(X(0)=0\)、\(X'(\ell)=0\) 给出
\[
 X_n(x)=\sin\dfrac{(n+1/2)\pi x}{\ell},
 \quad\lambda_n=\left(\dfrac{(n+1/2)\pi}{\ell}\right)^2,
 \quad n\geq0.
\]
这里最低频率不是零。决定零模态的条件是空间能量为零时哪些函数仍满足端点限制，而不是边界条件的名称中是否出现 Neumann。

\subsection{Robin 边界条件}
\label{b2:subsec:100104}

沿用\secref{b2:sec:0104}的外法向约定，区间两端的 Robin 条件写成
\begin{equation}\label{b2:eq:10-robin-flat}
 -X'(0)+h_0X(0)=0,
 \quad X'(\ell)+h_\ell X(\ell)=0,
 \quad h_0,h_\ell\in\mathbb R.
\end{equation}
于是
\[
 \lambda\|X\|_2^2
 =\|X'\|_2^2+h_0|X(0)|^2+h_\ell|X(\ell)|^2.
\]
当 \(h_0,h_\ell\geq0\) 时，谱非负；若其中至少一个严格为正，零特征函数只能为零。若二者均为零，就回到 Neumann 情形。负的边界系数会向系统输入能量，不能再由这个恒等式推出 \(\lambda\geq0\)。例如取 \(h_0=-a\)、\(h_\ell=a\)，\(a>0\)，则 \(X(x)=e^{-ax}\) 满足两个边界条件，特征值为 \(-a^2\)。

\ProseParagraphBreak

正频率仍可由一个标量方程求出。对 \(\lambda=k^2>0\)，左端条件使
\[
 X(x)=A\left(\cos(kx)+\dfrac{h_0}{k}\sin(kx)\right),
\]
代入右端条件得到
\begin{equation}\label{b2:eq:10-robin-characteristic}
 (h_0h_\ell-k^2)\sin(k\ell)
 +k(h_0+h_\ell)\cos(k\ell)=0.
\end{equation}
求根后还须归一化。零特征值则用 \(X=A+Bx\) 单独检查，其存在条件为 \(h_0+h_\ell+\ell h_0h_\ell=0\)。式\eqref{b2:eq:10-robin-characteristic}在 \(k=0\) 总形式消失，是前面乘以 \(k\) 造成的，不能把这个消失误认为零必是特征值。

\ProseParagraphBreak

显式三角函数仅适用于常系数模型。杆的密度、张力或热传导系数随位置变化时，特征函数一般不再等间距振荡。下一节不依赖显式求根，而从边界项、加权内积和解算子建立展开的共同基础。

% !TeX root = ../../Book2.tex
\section{Sturm–Liouville 问题}
\label{b2:sec:1002}

\subsection{特征值问题}
\label{b2:subsec:100201}

本节固定有限区间 \(I=(a,b)\)，实函数 \(p,w,q\in C([a,b])\)，并要求 \(p(x)>0\)、\(w(x)>0\) 在闭区间处处成立。连续性与紧性给出正的下界；\(q\) 可以变号。以下定义容许 \(p\) 只有连续性，不把 \((pu')'\) 强行拆成包含未定义的 \(p'\) 的表达式。

\begin{definition}\label{b2:def:10-sl-operator}
在每个端点 \(s\in\{a,b\}\) 固定一个条件：或为 \(u(s)=0\)，或为
\[
 \partial_{\nu,p}u(s)+h_su(s)=0,
 \quad h_s\in\mathbb R,
 \quad\partial_{\nu,p}u(a)=-p(a)u'(a),
 \quad\partial_{\nu,p}u(b)=p(b)u'(b).
\]
这些条件各只涉及一个端点，称为\term[fenli-bianjie-tiaojian]{分离边界条件}[separated boundary conditions]。记 \(R\) 为选择第二类表达式的端点集合，\(h_s=0\) 包含 Neumann 条件。令
\[
 H\coloneqq L^2(I,w\,\mathrm{d}x),
 \quad\langle u,v\rangle_H\coloneqq\int_a^b u\overline v\,w\,\mathrm{d}x,
\]
\[
 D(A)\coloneqq\{\,u\in H^1(I)\mid pu'\in H^1(I),\ u\text{ 满足已定端点条件}\,\},
 \quad Au\coloneqq\dfrac{-(pu')'+qu}{w}.
\]
方程 \(Au=\lambda u\)、\(u\in D(A)\setminus\{0\}\) 称为本章的\term[sturm-liouville-wenti]{Sturm--Liouville 特征值问题}[Sturm--Liouville eigenvalue problem]。
\end{definition}
\symidx{\(L^2(I,w\,\mathrm{d}x)\)：区间上的加权平方可积空间}

由第一卷命题 21.4.1，\(H^1(I)\) 的元素有唯一的绝对连续代表，并连续延至两端；\(pu'\in H^1\) 也有这种代表。因为 \(p\) 连续且不为零，\(u'=(pu')/p\) 有连续代表，故定义中的端点值和通量都有意义。这里 \(H\) 的内积仍对第一变量线性，\(w\) 同时出现在特征方程与正交关系中。若去掉这个权重，一般不会得到所需的对称性。

\ProseParagraphBreak

把 \(z=pu'\) 作为第二个未知量，齐次方程化成
\begin{equation}\label{b2:eq:10-first-order-system}
 u'=p^{-1}z,
 \quad z'=(q-\lambda w)u.
\end{equation}
给定一个端点的 \((u,z)\) 就唯一确定解。所需的常微分方程事实在这里只涉及连续系数线性系统：将积分方程反复代入，第 \(n\) 次差由 \(M^n|x-a|^n/n!\) 控制，迭代级数一致收敛，给出存在性；两个解之差反复使用同一估计则为零。故非零解不可能在同一点同时满足 \(u=z=0\)。这也解释了为何 \(p\) 在端点退化的 Legendre 型问题不属于本章范围。

\subsection{自伴性}
\label{b2:subsec:100202}

先处理积分恒等式，再检查算子定义域。对满足定义域正则性的 \(u,v\)，绝对连续函数的分部积分给出
\begin{equation}\label{b2:eq:10-lagrange-identity}
 \langle Au,v\rangle_H-\langle u,Av\rangle_H
 =\left[u\overline{pv'}-pu'\overline v\right]_a^b.
\end{equation}
这是一维的\term[lagrange-hengdengshi]{Lagrange 恒等式}[Lagrange identity]。同一分离边界条件使每个端点贡献单独消失：Dirichlet 端点有 \(u=v=0\)，Robin 端点则由实数 \(h_s\) 消去两项。因此 \(A\) 是对称的。无界算子的自伴性还要求伴随定义域与 \(D(A)\) 相同；此结论稍后由一个满射解算子证明，不能只由上式命名。

\ProseParagraphBreak

为找到一个可逆的移位，令 \(V\) 为 \(H^1(I)\) 中满足所有 Dirichlet 端点值为零的函数。Robin 条件不加入 \(V\) 的定义，而作为能量中的边界项出现。定义
\begin{equation}\label{b2:eq:10-energy-form}
 \mathfrak a(u,v)\coloneqq
 \int_a^b\left(pu'\overline{v'}+qu\overline v\right)\,\mathrm{d}x
 +\sum_{s\in R}h_su(s)\overline{v(s)},
 \quad u,v\in V.
\end{equation}
对 \(u\in D(A)\)、\(v\in V\)，一次分部积分给出 \(\mathfrak a(u,v)=\langle Au,v\rangle_H\)。这一步同时检验了两个端点的符号。

\begin{lemma}[移位能量估计]\label{b2:lem:10-form-coercive}
存在 \(\sigma>0\) 及 \(c,C>0\)，使
\[
 c\|u\|_{H^1(I)}^2
 \leq\mathfrak a_\sigma(u,u)
 \coloneqq\mathfrak a(u,u)+\sigma\|u\|_H^2
 \leq C\|u\|_{H^1(I)}^2
 \quad(u\in V).
\]
可以增加 \(\sigma\)，使 \(\mathfrak a_\sigma(u,u)\geq\|u\|_H^2\)。
\end{lemma}
\begin{proof}
对 \(0<r<b-a\)，由 \(u(a)=u(x)-\int_a^xu'\) 得
\[
 |u(a)|^2\leq2|u(x)|^2+2r\int_a^{a+r}|u'|^2.
\]
对 \(x\in(a,a+r)\) 平均，得到 \(|u(a)|^2\leq2r^{-1}\|u\|_2^2+2r\|u'\|_2^2\)。右端点有相同估计。给定 \(\varepsilon>0\)，取足够小的 \(r\)，便有
\[
 |u(a)|^2+|u(b)|^2
 \leq\varepsilon\|u'\|_2^2+C_\varepsilon\|u\|_2^2.
\]
用它控制所有负边界项，使导数项至少保留 \(\min p/2\)。负的 \(q\) 项由 \(\|q\|_\infty\|u\|_2^2\) 控制；再取 \(\sigma\) 足够大，由 \(\min w>0\) 吸收全部负的零阶项。上界则由系数有界性与固定 \(\varepsilon\) 的端点估计得到。
\end{proof}

特别地，\(A+\sigma I\) 的核为零。以下固定满足引理最后条件的 \(\sigma\)，并构造其逆。

\begin{proposition}[移位 Green 算子]\label{b2:prop:10-green-resolvent}
存在实连续对称核 \(G_\sigma(x,y)\)，使
\[
 (Kf)(x)\coloneqq\int_a^bG_\sigma(x,y)\cdot f(y)w(y)\,\mathrm{d}y
\]
对每个 \(f\in H\) 给出 \((A+\sigma I)u=f\) 的唯一解 \(u\in D(A)\)。作为 \(H\) 上的算子，\(K\) 有界、自伴、单射，而且 \(\langle Kf,f\rangle_H>0\) 对 \(f\ne0\) 成立。
\end{proposition}
\begin{proof}
取方程 \(-(pu')'+(q+\sigma w)u=0\) 的非零实解 \(\phi\)，使它满足左边界条件；另取满足右边界条件的非零实解 \(\psi\)。按端点数据构造这些解可用式\eqref{b2:eq:10-first-order-system}。量
\[
 W=p(x)\left(\phi'(x)\psi(x)-\phi(x)\psi'(x)\right)
\]
为常数：写成 \(z_\phi\psi-\phi z_\psi\) 后求导，两类项分别抵消。若 \(W=0\)，两组初值成比例，唯一性给出两解成比例，从而产生满足两端条件的非零齐次解，与移位能量估计矛盾。令
\[
 G_\sigma(x,y)=\dfrac1W
 \begin{cases}
 \phi(x)\psi(y),&x\leq y,\\
 \phi(y)\psi(x),&y\leq x.
 \end{cases}
\]
它连续、实且对称。固定 \(y\)，其 \(x\) 通量跨越 \(y\) 的跳跃为 \(-1\)，因而符号与负的二阶算子相符。

为避免仅凭形式跳跃断言，直接将 \(u=Kf\) 写成
\[
 u(x)=\dfrac1W\left(
 \psi(x)\int_a^x\phi(y)\cdot f(y)w(y)\,\mathrm{d}y
 +\phi(x)\int_x^b\psi(y)\cdot f(y)w(y)\,\mathrm{d}y\right).
\]
积分为绝对连续函数，求一阶导数时含 \(f(x)\) 的两项抵消。再次对通量求导得
\[
 (pu')'=(q+\sigma w)u-wf
\]
几乎处处成立。右侧属于 \(L^2\)，故 \(pu'\in H^1\)，而 \(u\in H^1\)。端点代入上式的积分表示及其一阶导数，得到相应条件。齐次差由能量估计为零，因此这是唯一解，且 \(K(A+\sigma I)u=u\) 对全部 \(u\in D(A)\) 成立。

连续核有界，Cauchy--Schwarz 不等式给出 \(K\) 的有界性；核对称与 Fubini 给出自伴性。若 \(Kf=0\)，应用 \(A+\sigma I\) 得 \(f=0\)。最后令 \(u=Kf\)，则
\[
 \langle Kf,f\rangle_H
 =\langle u,(A+\sigma I)u\rangle_H
 =\mathfrak a_\sigma(u,u)>0
\]
对非零 \(f\) 成立。
\end{proof}

这个积分核是移位边值问题的\term[green-hanshu]{Green 函数}[Green function]。实连续核的结论依赖有限规则区间；它不是高维 Green 函数在对角线也连续的断言。

\ProseParagraphBreak

对 Hilbert 空间中的稠密定义算子 \(T\)，若存在 \(g\) 使 \(\langle Tu,v\rangle=\langle u,g\rangle\) 对全部 \(u\in D(T)\) 成立，就令 \(v\in D(T^*)\)、\(T^*v=g\)；稠密性保证 \(g\) 唯一。这样得到的 \(T^*\) 称为\term[bansui-suanzi]{伴随算子}[adjoint operator]。若 \(D(T)=D(T^*)\) 且二者作用相同，则称 \(T\) 为\term[zi-ban-suanzi]{自伴算子}[self-adjoint operator]。全空间定义的有界算子是这个定义的特例。

\begin{theorem}[Sturm--Liouville 算子的自伴性]\label{b2:thm:10-self-adjoint}
\cref{b2:def:10-sl-operator}给出的 \(A\) 是 \(H\) 上稠密定义的闭自伴算子。
\end{theorem}
\begin{proof}
先证明稠密性。若 \(v\perp D(A)=\operatorname{ran}K\)，则 \(\langle Kf,v\rangle_H=\langle f,Kv\rangle_H=0\) 对一切 \(f\) 成立，所以 \(Kv=0\)，由单射性 \(v=0\)。闭性也由逆映射得到：若 \(u_j\to u\)、\(Au_j\to g\) 于 \(H\)，则 \(u_j=K(Au_j+\sigma u_j)\to K(g+\sigma u)\)，故 \(u=K(g+\sigma u)\in D(A)\) 且 \(Au=g\)。

对稠密定义的算子，\(v\in D(A^*)\) 指存在 \(g\in H\)，使 \(\langle Au,v\rangle_H=\langle u,g\rangle_H\) 对所有 \(u\in D(A)\) 成立，此时 \(A^*v=g\)。取 \(u=Kf\) 后得到
\[
 \langle f,v\rangle_H
 =\langle Kf,g+\sigma v\rangle_H
 =\langle f,K(g+\sigma v)\rangle_H.
\]
所以 \(v=K(g+\sigma v)\in D(A)\)，且 \(Av=g\)。这证明 \(A^*\subseteq A\)；Lagrange 恒等式给出反向包含，故两者定义域和作用均相同。
\end{proof}

\subsection{正交性}
\label{b2:subsec:100203}

\begin{proposition}\label{b2:prop:10-real-orthogonal-simple}
本章的特征值均为实数，不同特征值的特征函数在 \(H\) 中正交，每个特征空间为一维。
\end{proposition}
\begin{proof}
若 \(Au=\lambda u\)，则 \(\lambda\|u\|_H^2=\mathfrak a(u,u)\) 为实数。若再有 \(Av=\mu v\)，自伴性给出 \((\lambda-\mu)\langle u,v\rangle_H=0\)。最后，满足左分离条件的初值 \((u(a),p(a)u'(a))\) 构成二维空间中的一条直线。同一 \(\lambda\) 的任何两个非零解在左端初值成比例，常微分方程唯一性使两解处处成比例。
\end{proof}

单重性依赖分离条件。周期条件耦合两个端点，正弦与余弦可以属于同一特征值；它们的二维特征空间并不违反本命题。正交性本身更为一般，只需所选定义域使边界配对消失。

\subsection{特征值离散性}
\label{b2:subsec:100204}

能量下界已经排除了充分负的特征值，还须排除有限处的无穷聚集。若 \(\lambda_j\) 为有界的不同特征值，相应单位特征函数 \(e_j\) 两两正交，且
\[
 \mathfrak a_\sigma(e_j,e_j)=\lambda_j+\sigma
\]
一致有界。移位估计使它们在 \(H^1(I)\) 中有界。一维积分表示给出 \(|e_j(x)-e_j(y)|\leq C|x-y|^{1/2}\)，且端点估计使函数一致有界。对每个正整数 \(m\) 选网格宽度小于 \(1/m\) 的有限网格，依次取子列，使第一个网格上的全部函数值收敛，再使第二个网格上的全部值收敛，如此继续。对角子列在每个固定网格上收敛；任意点可移到距它至多 \(1/m\) 的网格点，误差至多 \(Cm^{-1/2}\)。先取 \(m\) 大，再比较子列中足够靠后的两项，就得到一致 Cauchy 性。因此存在一致收敛子列，特别在 \(H\) 中收敛；这与任意两项距离 \(\sqrt2\) 矛盾。下一节用 \(K\) 的紧性把同一机制提升为完备性定理，同时证明特征值确有无穷多个。

% !TeX root = ../../Book2.tex
\section{完备性}
\label{b2:sec:1003}

\subsection{紧算子方法}
\label{b2:subsec:100301}

算子 \(A\) 对高频求导，通常无界；其移位逆 \(K\) 则是具有紧性的积分算子。若一个有界线性算子把每个有界列送到有收敛子列的像列，则称它为\term[jin-suanzi]{紧算子}[compact operator]。这里使用第一卷定理 8.5.5 的紧自伴谱定理：紧自伴算子的非零特征向量张成核空间的正交补，每个非零特征值重数有限，唯一可能的聚点是零。应用前须核对紧性及核，而不能只从自伴性猜测完备性。

\begin{theorem}[Sturm--Liouville 展开]\label{b2:thm:10-sl-spectral}
在\cref{b2:def:10-sl-operator}的假设下，存在一组实的标准正交特征函数 \((e_n)_{n\geq1}\)，构成 \(H\) 的基。相应特征值可按严格递增次序排列为
\[
 -\sigma<\lambda_1<\lambda_2<\cdots,
 \quad\lambda_n\longrightarrow+\infty.
\]
特别地，负特征值只有有限多个。每个 \(f\in H\) 都有
\[
 f=\sum_{n=1}^{\infty}f_ne_n,
 \quad f_n=\langle f,e_n\rangle_H,
 \quad\|f\|_H^2=\sum_{n=1}^{\infty}|f_n|^2,
\]
其中级数在 \(H\) 中收敛。
\end{theorem}
\begin{proof}
将 \([a,b]^2\) 分成有限个小矩形，在每个矩形内以 \(G_\sigma\) 的一个取值近似它。连续性给出阶梯核 \(G_j\to G_\sigma\) 一致收敛。对应算子 \(K_j\) 的像落在有限个区间示性函数的张成空间内，因而有限秩。记 \(m_w=\int_a^bw\)，则
\[
 \|(K-K_j)f\|_H
 \leq m_w\|G_\sigma-G_j\|_\infty\cdot\|f\|_H.
\]
故 \(K_j\to K\) 于算子范数，\(K\) 紧。第一卷紧自伴谱定理与\cref{b2:prop:10-green-resolvent}给出正特征值 \(\mu_n\) 及其标准正交特征向量；因 \(\ker K=0\)，这些向量在全部 \(H\) 中完备。这个列必须无穷，因为 \(H\) 无穷维，而有限秩算子不能在无穷维空间上单射。因此 \(\mu_n\to0\)。

若 \(Ke_n=\mu_ne_n\)，则 \(e_n\in\operatorname{ran}K=D(A)\)，且 \(Ae_n=(\mu_n^{-1}-\sigma)e_n\)。反过来，任一 \(A\) 的特征向量都是 \(K\) 的特征向量。上一节的单重性使 \(\mu_n\) 无重复，并可取实特征函数；实系数常微分方程的左端初值先取为实数即可。按 \(\mu_n\) 递减排序，得到所述严格递增的 \(\lambda_n\)。完备性给出展开及 Parseval 等式。
\end{proof}

这里的两个谱极限方向相反：\(K\) 的小特征值对应 \(A\) 的大特征值。零始终是 \(K\) 的谱极限，却不是它的特征值；\(A\) 是否有零特征值，则取决于原方程及边界条件。移位只是为证明选择一个可逆点，并没有把物理模型中的零频率删去。

\ProseParagraphBreak

有限区间的假设不能仅从符号上删掉。考虑实线上 \(-u''+u=f\)，其平方可积解算子为与 \(k(x)=e^{-|x|}/2\) 卷积。分布求导给出 \(-k''+k=\delta_0\)，而 \(\widehat k(\xi)=(1+4\pi^2\xi^2)^{-1}\)，所以这是有界单射的解算子。然而它不紧：固定非零 \(f\in C_c^\infty\)，记 \(g=k*f\ne0\)。平移后的 \(f(\,\cdot-jR)\) 范数不变，相应的解为 \(g(\,\cdot-jR)\)。平方可积函数的远平移内积趋零，这可先把 \(g\) 在大紧区间外的平方尾截掉，再使用两个紧支集的分离证明。选足够稀疏的平移序列，相应解之间的距离便有同一个正下界，不可能有收敛子列。

\ProseParagraphBreak

这个反例中，局部积分仍然提高正则性，但函数可以向无穷远整体移动，阻止了全空间紧性。本章连续核在紧正方形上的有限秩逼近同时利用了局部平滑和空间范围有界；后一项在实线模型中消失。因而全空间的 Fourier 积分不能无条件改写成一列平方可积特征函数的展开。

\subsection{特征函数展开}
\label{b2:subsec:100302}

为了使用展开求解方程，需要知道哪些系数列还可以求导。记 \(r_n=\lambda_n+\sigma\geq1\)，并令 \(P_Nf=\sum_{n=1}^Nf_ne_n\)。以下两个刻画区分算子定义域与能量空间。

\begin{proposition}[定义域和能量的系数刻画]\label{b2:prop:10-spectral-domains}
对 \(f\in H\)，有
\[
 f\in D(A)\quad\Longleftrightarrow\quad
 \sum_n r_n^2|f_n|^2<\infty,
\]
此时 \(Af=\sum_n\lambda_nf_ne_n\) 在 \(H\) 中成立。又有
\begin{equation}\label{b2:eq:10-form-parseval}
 f\in V\quad\Longleftrightarrow\quad\sum_n r_n|f_n|^2<\infty,
 \quad\mathfrak a_\sigma(f,f)=\sum_n r_n|f_n|^2.
\end{equation}
对应截断分别在定义域范数 \(\|f\|_H+\|Af\|_H\) 与 \(H^1(I)\) 范数中收敛。
\end{proposition}
\begin{proof}
若 \(f\in D(A)\)，令 \(g=(A+\sigma I)f\)，则 \(f=Kg\)，系数关系为 \(g_n=r_nf_n\)。反过来，若 \((r_nf_n)\in\ell^2\)，它给出某个 \(g\in H\)，而 \(Kg\) 与 \(f\) 的全部系数相同，故 \(f=Kg\in D(A)\)。这也证明定义域范数中的截断收敛。

\(V\) 是 \(H^1\) 的闭子空间，因为端点取值连续；由\cref{b2:lem:10-form-coercive}，\(\mathfrak a_\sigma\) 是它的一个等价 Hilbert 内积。函数 \(e_n/\sqrt{r_n}\) 在这个内积下标准正交。若 \(v\in V\) 与它们全部正交，分部积分给出
\[
 0=\mathfrak a_\sigma(v,e_n)=r_n\langle v,e_n\rangle_H,
\]
而 \(H\) 中的完备性推出 \(v=0\)。所以它们也是 \(V\) 的标准正交基，\(V\) 中的 Fourier 系数为 \(\sqrt{r_n}v_n\)，由此得式\eqref{b2:eq:10-form-parseval}的正向及等式。反向中若加权和有限，\(P_Nf\) 在 \(V\) 中为 Cauchy 列，其 \(V\) 极限在 \(H\) 中只能为 \(f\)，故 \(f\in V\)。
\end{proof}

系数刻画也给出非齐次边值问题的完整可解条件。若实数 \(z\) 不等于任何 \(\lambda_n\)，则
\[
 u=\sum_n\dfrac{f_n}{\lambda_n-z}e_n
\]
属于 \(D(A)\)，并且 \((A-zI)u=f\)。因为 \(\lambda_n\to\infty\)，比值 \(r_n/|\lambda_n-z|\) 一致有界，定义域条件正由这个界保证。若 \(z=\lambda_k\)，可解当且仅当 \(f_k=0\)；满足这个条件时去掉第 \(k\) 项所得级数给出一个解，全部解还可加上任意倍 \(e_k\)。这里的自由常数来自齐次核，并非级数方法漏掉的解。

\ProseParagraphBreak

例如 Neumann 问题 \(-u''=f\) 的可解条件为 \(\int_0^\ell f=0\)，因为零特征函数为常数。将方程直接积分也得到这一条件；谱方法进一步给出它的充分性和除常数外的唯一性。取零均值解相当于在核的正交补中选取唯一代表。

\subsection{Fourier 级数作为特殊情形}
\label{b2:subsec:100303}

对 \(p=w=1,q=0\)，\cref{b2:thm:10-sl-spectral}与显式计算给出正弦基、余弦基及混合边界基的完备性。正弦展开也可以把区间函数作奇延拓后由周期 Fourier 理论得到，余弦展开则对应偶延拓。两条路线各有用途：周期化便于使用卷积核和求和公式，Green 算子路线则能在 \(p,w,q\) 改变后继续成立。

\ProseParagraphBreak

能量刻画把熟悉的导数公式变成精确命题。对 \(f\in H_0^1(0,\ell)\)，其正弦系数满足
\[
 \int_0^\ell|f'|^2\,\mathrm{d}x
 =\sum_{n\geq1}\left(\dfrac{n\pi}{\ell}\right)^2|f_n|^2.
\]
因而
\[
 \|f\|_2\leq\dfrac\ell\pi\|f'\|_2,
\]
等号恰在第一正弦模态及零函数处成立。Neumann 情形对零均值的 \(H^1\) 函数有同一最优常数，但不加零均值条件时常数函数使不等式失败。这些结论来自同一能量恒等式；边界值与均值限制决定了最低允许频率。

\ProseParagraphBreak

一般系数问题仍有 \(H\) 与 \(V\) 中的截断收敛，却不能仅由本章谱定理断言任意连续函数的原始特征函数级数处处收敛。正交完备、点态收敛与一致收敛是不同层次。下一节以范数极限解释演化方程，正是为避免把第一种结论误当成后两种。


% !TeX root = ../../Book2.tex
\section{广义解与逼近}
\label{b2:sec:1004}

\subsection{形式级数解}
\label{b2:subsec:100401}

本节在 \(\mathfrak a(u,u)\geq0\) 的附加假设下讨论演化，故 \(\lambda_n\geq0\)。例如 \(q\geq0\)、所有 Robin 系数非负时即满足此条件。空间仍取上一节的 \(H,V,A\)，时间系数吸收进 \(A\)。对方程
\[
 u_{tt}+Au=0,
 \quad u(0)=f,
 \quad u_t(0)=g,
\]
逐模态的候选为
\begin{equation}\label{b2:eq:10-wave-series}
 u(t)=\sum_n\left(f_n\cos(\sqrt{\lambda_n}t)
 +g_n s_n(t)\right)e_n,
 \quad s_n(t)=
 \begin{cases}
 \sin(\sqrt{\lambda_n}t)/\sqrt{\lambda_n},&\lambda_n>0,\\
 t,&\lambda_n=0.
 \end{cases}
\end{equation}
而 \(u_t+Au=0\)、\(u(0)=f\) 的候选为 \(\sum_ne^{-\lambda_nt}f_ne_n\)。每个有限截断都可逐项求导，它满足方程及截断后的数据。现在的问题是：这些有限解以什么方式趋近一个满足原数据的解？

\subsection{不可逐项经典求导时的处理}
\label{b2:subsec:100402}

取正弦系，若 \(f_n=n^{-2}\)、\(g_n=0\)，则 \(\sum_n\lambda_n|f_n|^2<\infty\)，但 \(\sum_n\lambda_n^2|f_n|^2=\infty\)。这给出有限能量初位移，却不属于 \(D(A)\)。因此波方程在初始时刻不必有 \(H\) 值二阶时间导数；把 \(-\lambda_nf_n\) 的级数硬作 \(L^2\) 函数会失去收敛。合理的做法是把空间导数移到检验函数上。

\begin{definition}\label{b2:def:10-energy-wave}
设 \(f\in V\)、\(g\in H\)。若 \(u\in C([0,T];V)\)、\(u_t\in C([0,T];H)\)，初值分别为 \(f,g\)，并且对每个 \(v\in V\)、\(\eta\in C_c^\infty(0,T)\) 有
\[
 \int_0^T\left(-\langle u_t(t),v\rangle_H\eta'(t)
 +\mathfrak a(u(t),v)\eta(t)\right)\,\mathrm{d}t=0,
\]
则称 \(u\) 为波动问题的能量解。此处将\chapref{b2:ch:09}的能量解概念用于带分离边界的加权区间。
\end{definition}

条件 \(u(t)\in V\) 保存 Dirichlet 端点值。Robin 通量一般不能对能量解逐点求值，其边界信息已包含在 \(\mathfrak a\) 中。正则性提高到 \(D(A)\) 时，分部积分又恢复原来的通量条件。因此这种定义保留了边界约束的作用，并没有要求所有弱解都有经典边界导数。

\begin{theorem}[波动问题的能量极限]\label{b2:thm:10-wave-energy-limit}
对每个 \(f\in V\)、\(g\in H\)，式\eqref{b2:eq:10-wave-series}在每个有限时间区间上给出唯一能量解。它的有限截断在 \(C([0,T];V)\) 中收敛，时间导数在 \(C([0,T];H)\) 中收敛，并满足
\begin{equation}\label{b2:eq:10-wave-energy}
 \dfrac12\|u_t(t)\|_H^2+\dfrac12\mathfrak a(u(t),u(t))
 =\dfrac12\|g\|_H^2+\dfrac12\mathfrak a(f,f).
\end{equation}
若再有 \(f\in D(A)\)、\(g\in V\)，则 \(u\in C([0,T];D(A))\cap C^1([0,T];V)\cap C^2([0,T];H)\)。
\end{theorem}
\begin{proof}
每个标量系数 \(c_n\) 满足 \(c_n''+\lambda_nc_n=0\)，故
\[
 |c_n'(t)|^2+\lambda_n|c_n(t)|^2
 =|g_n|^2+\lambda_n|f_n|^2.
\]
截断之差对任意尾部满足相同等式。初数据条件使右侧尾和趋零，一致控制 \(H\) 值速度和 \(\mathfrak a\) 能量。另由
\[
 \|u_N(t)-u_M(t)\|_H
 \leq\|P_Nf-P_Mf\|_H
 +T\sup_{s\leq T}\|u_N'(s)-u_M'(s)\|_H
\]
补上可能存在的零模态。结合 \(\mathfrak a_\sigma\) 与 \(H^1\) 的等价性，得到所称两种一致收敛，极限的时间导数由积分恒等式 \(u_N(t)=P_Nf+\int_0^tu_N'(s)\,\mathrm{d}s\) 传递。

有限截断满足定义中的检验恒等式；\(\mathfrak a\) 在 \(V\times V\) 连续，所以可以传递到极限。能量等式也由范数收敛传递。对任一能量解，取 \(v=e_n\)，得到系数 \(c_n\) 的标量微分方程及初值；其分布二阶导数为连续函数 \(-\lambda_nc_n\)，积分两次后即为经典方程，故系数唯一。完备性再给出解的唯一性。

若初值更正则，将刚才的标量恒等式乘以 \(\lambda_n\)，加上原等式和 \(H\) 分量控制，便得到 \(\sum r_n^2|c_n(t)|^2\)、\(\sum r_n|c_n'(t)|^2\) 的统一可和尾界。由\cref{b2:prop:10-spectral-domains}得前三个空间的连续性，且 \(u_{tt}=-Au\) 于 \(H\) 中连续。
\end{proof}

\subsection{能量极限解释}
\label{b2:subsec:100403}

热方程允许更粗的初值，因为每个正频率都有指数衰减。为了将此结论用于\chapref{b2:ch:12}，先给出不依赖逐点求导的版本。

\begin{proposition}[热方程的平方能量极限]\label{b2:prop:10-heat-energy-limit}
对 \(f\in H\)，级数
\[
 u(t)=\sum_ne^{-\lambda_nt}f_ne_n
\]
定义 \(C([0,\infty);H)\) 函数，\(u(0)=f\)，且对 \(t>0\) 属于 \(D(A)\)，在 \(H\) 中满足 \(u_t=-Au\)。对 \(T>0\)，有
\begin{equation}\label{b2:eq:10-heat-energy}
 \|u(T)\|_H^2+2\int_0^T\mathfrak a(u(t),u(t))\,\mathrm{d}t
 =\|f\|_H^2.
\end{equation}
有限截断在 \(C([0,T];H)\) 以及 \(L^2((0,T);V)\) 中收敛。
\end{proposition}
\begin{proof}
对全部 \(t\geq0\)，平方尾和不超过 \(\sum_{n>N}|f_n|^2\)，故有限截断一致收敛于 \(H\) 值连续函数并取得初值。在 \(t\geq\delta>0\) 上，\(r_n^2e^{-2\lambda_nt}\) 一致有界；更高有限次幂也有同一性质。因此定义域刻画允许在 \(H\) 中逐项求一次时间导数，得到 \(u_t=-Au\)。

逐模态积分恒等式为
\[
 e^{-2\lambda_nT}|f_n|^2
 +2\int_0^T\lambda_ne^{-2\lambda_nt}|f_n|^2\,\mathrm{d}t
 =|f_n|^2,
\]
在 \(\lambda_n=0\) 时也成立。由非负项求和与单调收敛，得到式\eqref{b2:eq:10-heat-energy}。对尾部应用同一式，给出积分能量尾趋零；再加 \(\sigma T\sum_{n>N}|f_n|^2\)，由能量范数等价得到 \(L^2((0,T);V)\) 收敛。
\end{proof}

这个极限也满足弱方程。对 \(v\in V\)、\(\eta\in C_c^\infty(0,T)\)，由有限解再过极限有
\[
 \int_0^T\left(-\langle u(t),v\rangle_H\eta'(t)
 +\mathfrak a(u(t),v)\eta(t)\right)\,\mathrm{d}t=0.
\]
在 \(C([0,T];H)\cap L^2((0,T);V)\) 的同初值解类中，取 \(v=e_n\) 就得到 \(c_n'+\lambda_nc_n=0\)，从而唯一。这里不需要先给全部弱解建立可逐点微分的能量等式。

\ProseParagraphBreak

若原算子有负特征值，谱展开仍成立，但波动中的相应三角函数变为双曲函数，热方程相应项按 \(e^{|\lambda_n|t}\) 增长。只有有限个负模态，故有限时间内的解理论仍可逐项处理；收缩性、非负能量和趋向平衡等结论却必须重新检查。本章单独陈述非负情形的结论，后续讨论扩散与振动模型时将用到这些结果。

% 本章小结（不计入编号）
\begin{chaptersummary}
空间边界条件决定允许的模态，时间方程再决定它们振动、衰减或增长。常系数时可以显式计算；一般连续系数时，移位 Green 逆把微分问题变成一个紧自伴算子问题。逆算子的核为零，因而特征函数在整个加权平方可积空间中完备。

\ProseParagraphBreak

完备性只先给出平方范数中的展开。算子定义域需要两次频率权，能量空间只需一次；这一区别决定能否经典求导，以及应把边界条件理解为端点值还是弱能量项。有限模态解的统一估计给出波动与扩散的能量极限，后续讨论具体方程时可以直接使用这些结论。
\end{chaptersummary}


\BookExercises

前四题训练边界、权重和变换；后四题考察稳定性、参数变化与演化。各题均应区分零模态和正模态。

% 自拟：混合边界的半整数频率及最优常数，增加与双Dirichlet的比较。
\begin{exercise}\label{b2:ex:10-mixed-best-constant}
在 \((0,\ell)\) 上取 \(u(0)=0\)、\(u'(\ell)=0\)。求 \(-u''\) 的单位特征函数，并证明对每个 \(u\in H^1(0,\ell)\)、\(u(0)=0\)，有最优不等式 \(\|u\|_2\leq(2\ell/\pi)\|u'\|_2\)。能否再要求 \(u'(\ell)=0\) 来定义这里的全部能量函数？
\end{exercise}
\begin{solution}
特征函数为 \(\sqrt{2/\ell}\sin((n+1/2)\pi x/\ell)\)，\(n\geq0\)，特征值是相应频率平方。能量展开给出 \(\|u'\|_2^2\geq(\pi/(2\ell))^2\|u\|_2^2\)，第一模态取等号。一般 \(H^1\) 函数只有函数本身的端点迹，导数仅属 \(L^2\)，没有定义导数端点值。因此自然 Neumann 条件属于算子定义域，而不应额外写进能量空间。
\end{solution}

% 自拟：零频根检查与非正边界的显式负模态。
\begin{exercise}\label{b2:ex:10-robin-zero}
对式\eqref{b2:eq:10-robin-flat}，完整分类零特征值存在的参数，并写出其特征函数。再取 \(h_0=-a,h_\ell=a\)，\(a>0\)，求全部特征值。
\end{exercise}
\begin{solution}
零频通解 \(A+Bx\) 的左条件给出 \(B=h_0A\)，非零解必有 \(A\ne0\)。右条件等价于 \(h_0+h_\ell+\ell h_0h_\ell=0\)，特征函数为 \(1+h_0x\)。对指定参数，此式变为 \(-\ell a^2=0\)，所以没有零频。正频方程化成 \(-(a^2+k^2)\sin(k\ell)=0\)，给出 \((n\pi/\ell)^2\)，\(n\geq1\)。负频 \(\lambda=-s^2\) 的同一双曲计算给出 \((s^2-a^2)\sinh(s\ell)=0\)，故唯一负特征值为 \(-a^2\)，特征函数 \(e^{-ax}\)。
\end{solution}

% 自拟：变量代换同时改变主系数与权重，不能漏掉加权空间。
\begin{exercise}\label{b2:ex:10-weighted-coordinate}
设 \(\rho\in C([a,b])\) 且严格为正，取 \(p=\rho^{-1}\)、\(w=\rho\)、\(q=0\)，两端为 Dirichlet。用 \(y=\int_a^x\rho(s)\,\mathrm{d}s\) 把问题化为常系数问题，求全部特征值与单位特征函数。
\end{exercise}
\begin{solution}
记 \(L=\int_a^b\rho\)，\(u(x)=v(y)\)。则 \(u'=\rho v_y\)、\(pu'=v_y\)，而 \(-(pu')'=-\rho v_{yy}\)，故特征方程为 \(-v_{yy}=\lambda v\)。又 \(\int_a^b|u|^2\rho\,\mathrm{d}x=\int_0^L|v|^2\,\mathrm{d}y\)，所以变换为等距同构。于是 \(\lambda_n=(n\pi/L)^2\)，\(e_n(x)=\sqrt{2/L}\sin((n\pi/L)\int_a^x\rho)\)。使用无权 \(L^2\) 归一化会破坏这个直接对应。
\end{solution}

% 自拟：周期边界展示分离假设对单重性的必要性。
\begin{exercise}\label{b2:ex:10-periodic-multiplicity}
在 \((0,\ell)\) 上改取周期条件 \(u(0)=u(\ell)\)、\(u'(0)=u'(\ell)\)。直接证明边界配对消失，求 \(-u''\) 的特征值与重数，并说明为何不违背\cref{b2:prop:10-real-orthogonal-simple}。
\end{exercise}
\begin{solution}
Lagrange 恒等式的右端两个端点值相同，作差为零。能量非负，零特征空间是常数。正频通解的系数向量经过长度 \(\ell\) 的旋转后必须不变，故 \(k\ell=2n\pi\)。特征值 \((2n\pi/\ell)^2\)、\(n\geq1\)，各有正弦、余弦两个线性无关特征函数。周期条件耦合两端，不能用一条左端初值直线描述，故不满足分离假设。
\end{solution}

% 自拟：逼近共振时的误差放大，连接闭算子与求解稳定性。
\begin{exercise}\label{b2:ex:10-resonance-norm}
设 \(z\in\mathbb R\setminus\{\lambda_n\}\)。证明
\[
 \|(A-zI)^{-1}\|_{H\to H}
 =\dfrac1{\operatorname{dist}(z,\{\lambda_n\})}.
\]
若 \(f_k\ne0\)，描述 \(z\to\lambda_k\) 时解的主部；若 \(f_k=0\)，说明在 \(e_k^\perp\) 中极限为何存在。
\end{exercise}
\begin{solution}
系数展开给出范数上界；取到距离最小的单位特征函数得到反向不等式。距离的最小值存在，因为特征值仅向正无穷发散。解可拆为 \(f_k(\lambda_k-z)^{-1}e_k\) 与其余级数。其余分母在 \(z\) 的充分小邻域内一致远离零，且 \(r_n/|\lambda_n-z|\) 有统一界，故余项甚至在定义域范数中趋于删去第 \(k\) 项的解。主部是否出现恰由 \(f_k\) 决定。
\end{solution}

% 自拟：从已证能量展开导出极小极大公式，再比较势函数。
\begin{exercise}\label{b2:ex:10-potential-comparison}
证明
\[
 \lambda_n=\min_{\substack{E\subset V\\\dim E=n}}
 \max_{u\in E\setminus\{0\}}\dfrac{\mathfrak a(u,u)}{\|u\|_H^2}.
\]
由此证明固定 \(p,w\) 和边界条件，若 \(q_1\leq q_2\)，则 \(\lambda_n(q_1)\leq\lambda_n(q_2)\)；若 \(|q_1-q_2|\leq\varepsilon w\)，则两个特征值之差的绝对值不超过 \(\varepsilon\)。
\end{exercise}
\begin{solution}
每个 \(n\) 维子空间都含一个非零向量与 \(e_1,\ldots,e_{n-1}\) 正交。能量展开使其商不小于 \(\lambda_n\)。取 \(E=\operatorname{span}(e_1,\ldots,e_n)\) 时商的最大值等于 \(\lambda_n\)，故极值公式成立。对同一个 \(u\)，势函数增大使能量商增大；逐次取最大与最小保持不等号。最后的假设使两个商之差至多为 \(\varepsilon\)，双向比较即得结论。
\end{solution}

% 自拟：有限能量不足以保证二阶空间导数平方可积。
\begin{exercise}\label{b2:ex:10-energy-not-strong}
对 \((0,\pi)\) 的 Dirichlet 问题，令 \(f=\sum_{n\geq1}n^{-2}e_n\)。证明 \(f\in H_0^1\setminus D(A)\)。波动方程取初位移 \(f\)、初速度零，证明它的初始能量有限，但 \(u_{tt}(0)\) 不可能是 \(H\) 中的强二阶导数。
\end{exercise}
\begin{solution}
此时 \(\lambda_n=n^2\)。能量和为 \(\sum n^{-2}<\infty\)，定义域平方和却含 \(\sum1\)，故前一结论成立。若强二阶导数存在，对连续线性泛函 \(v\mapsto\langle v,e_n\rangle\) 求导，它的第 \(n\) 个系数必须为 \(-n^2n^{-2}=-1\)。这个系数列不属于 \(\ell^2\)，与 \(H\) 中的 Parseval 等式矛盾。能量解的弱二阶导数仍然存在。
\end{solution}

% 自拟：外力激发单一模态，区分稳定振荡与共振增长。
\begin{exercise}\label{b2:ex:10-forced-mode}
设 \(\lambda_k>0\)。求零初值问题 \(u_{tt}+Au=\cos(\omega t)e_k\) 的解，其中 \(\omega\geq0\)。分析 \(\omega=\sqrt{\lambda_k}\) 与其他情形，并说明有限区间的离散谱如何决定共振频率。
\end{exercise}
\begin{solution}
只有第 \(k\) 个系数非零。若 \(\omega^2\ne\lambda_k\)，解为
\[
 u(t)=\dfrac{\cos(\omega t)-\cos(\sqrt{\lambda_k}t)}{\lambda_k-\omega^2}e_k.
\]
分子及一阶导数在零时为零，直接求导检验方程。令 \(\omega\to\sqrt{\lambda_k}\)，或重新积分标量方程，得到共振解
\[
 u(t)=\dfrac{t\sin(\sqrt{\lambda_k}t)}{2\sqrt{\lambda_k}}e_k.
\]
此时振幅线性增长；其他频率的响应有界。能量增长来自外力，不能套用齐次问题的守恒式。
\end{solution}
